\scanpage{5}

On dit alors que les formes $\varphi$ et $\psi$ \deleted{(satisfaisant aux conditions de Weil)} sont \uline{compatibles}.

4) La compatibilité \deleted{des} formes de Weil est une relation \deleted{d'équivalence} réflexive et symétrique. \deleted{\unsure{E|De}} \inserted{\deleted{De façon plus précise, si $\varphi$ est une forme de}} $\varphi,\varphi'$ sur $M$, $\psi,\psi'$ sur $N$ et \deleted{$\chi$ \unsure{sont deux des formes de} Weil, si $\varphi$ est compatible à $\psi$, $\psi$ à $\chi$, et s'il existe une forme} \inserted{si $\varphi,\varphi'$ compatibles, $\psi,\psi'$ compatibles,} alors $\varphi$ et $\psi$ sont compatibles ssi $\varphi'$ et $\psi'$ le sont. Pour le voir, il faut expliciter :

% slanted
\margin{\deleted{\unsure{$\varphi,\psi$, $\varphi$ est} \unsure{compatible} \unsure{ssi $\varphi$ et $\psi$}} relation pas transitive, car si $M,N$ sont tels que $\operatorname{Hom}(M,N)=\operatorname{Hom}(N,M)=0$, $M\ne0$, alors pour deux \deleted{formes $\varphi,\unsure{\psi}$} \inserted{de Weil} $\varphi$ sur $M$, $\psi$ sur $N$, on a $\varphi$ comp. à $\psi$, $\psi$ comp. avec $-\varphi$ mais pas $\varphi$ avec $-\varphi$.}

5) \deleted{Soient} $\varphi$ et $\psi$ deux formes $\varepsilon$-symétriques $M\otimes M\to L$, \deleted{\unsure{on}} $u\mapsto u^*$ de $\operatorname{End}(M)$ associée à $\varphi$, et $\psi$ définie par $\alpha\in\operatorname{End}(M)$, $\psi(x,y)=\varphi(x,\alpha y)$, i.e. $\widetilde\psi=\widetilde\varphi\circ\alpha$. Alors $u^{\operatorname{ad}(\psi)}$ défini par $\psi$ est \deleted{\unsure{donné par}} \inserted{égal à} $\alpha^{-1}u^*\alpha$, \deleted{donc $\varphi$ $\psi$ est compatible avec $\varphi$ ssi $\operatorname{Tr}$} et $u^{\operatorname{ad}(\varphi,\psi)}$ \deleted{\unsure{défini|défini par $\psi$}} égal à $u^*\alpha$. \deleted{Donc}

\[
\begin{aligned}
&\operatorname{Tr}u\,u^{\operatorname{ad}(\varphi,\psi)}>0 &&\text{pour tt }u\ne0\\
&\hspace{2.7em}\Updownarrow\\
&\operatorname{Tr}uu^*\alpha>0 &&\text{pour tt }u\ne0\\
&\hspace{2.7em}\rotatebox{90}{$=$}\\
&\hspace{1.6em}\operatorname{Tr}u^*\alpha u
\end{aligned}
\]

De plus, cette condition implique que \deleted{pour tt} \inserted{$\psi$ est de Weil}, i.e.

\[
\begin{aligned}
&\operatorname{Tr}u\,u^{\operatorname{ad}(\psi,\psi)}>0 &&\text{pour }u\ne0\\
\text{i.e.}\quad &\operatorname{Tr}u\alpha^{-1}u^*\alpha>0 &&\text{pour }u\ne0
\end{aligned}
\]

Donc les formes de Weil sur $M$ compatibles à $\varphi$ correspondent aux \deleted{\unsure{homs}} $\alpha\in\operatorname{End}(M)$ tels que $\operatorname{Tr}u^*u\alpha>0$ pour $u\ne0$. \inserted{\deleted{\unsure{donc}} Donc, pour les formes \inserted{de Weil} sur $M$, la relation de compatibilité est une relation d'équivalence.}

% slanted
\margin{Conséquence : si $\varphi$ et $\psi$ sont de Weil compatibles sur un $=$ $M$, $\varphi\unsure{+\!|\circ|\oplus}\psi$ est de Weil et compatible avec $\varphi$ et $\psi$}

6) Soient \deleted{\unsure{$\varphi_i$}} $(M_i,\varphi_i)$, \deleted{des} avec $\varphi_i:M_i\otimes M_i\to L$, $\varphi_i$ $\varepsilon$-symétriques de Weil ($\varepsilon=\pm1$ fixé). Supposons que pour $i,j$, $\varphi_i$ et $\varphi_j$ compatibles. Alors

% archive sheet number: 4
